%% ── Ⅴ. CONCLUSION ───────────────────────────────────────
\chapter{Conclusion}

The conclusion generally consists of a single chapter and includes discussions of the research results, conclusions, and suggestions for future work.

%% ─────────────────────────────────────────────────────────
\section{Discussion}

Function of the discussion:
\begin{itemize}[leftmargin=3em, itemsep=2pt]
  \item To provide answers to the research questions presented in the introduction
  \item To explain the results supporting the answers to the research questions and describe how these answers appropriately fit within the current body of knowledge
\end{itemize}

Studies on the impact and structural behavior of nuclear containment structures have been conducted using various methodologies.
Jung analyzed the nonlinear dynamic behavior of CANDU-type nuclear containment structures subjected to aircraft impact loads \cite{jung2002},
while Choi performed a local impact analysis of prestressed nuclear containment walls considering the effect of prestressing \cite{choi2014}.
In addition, Kyratsis investigated children's interactions and the co-construction of peer culture \cite{kyratsis2004}.
In the field of anxiety and depression treatment, Wells proposed a metacognitive therapy approach \cite{wells2009a,wells2009b,wells2010}.
Milnes analyzed the possibility of reducing adolescent depression using generative instruction methods \cite{milnes1998}.
Furthermore, explanations regarding website citation methods can be found in Caulfield’s APA style guide \cite{caulfield2020}.